\section{Cohen-Macaulay rings}
\label{section-CM-ring}

\noindent
Most of the results of this section are special cases of the results
in Section \ref{section-CM}.

\begin{definition}
\label{definition-local-ring-CM}
A Noetherian local ring $R$ is called {\it Cohen-Macaulay}
if it is Cohen-Macaulay as a module over itself.
\end{definition}

\noindent
Note that this is equivalent to requiring the existence
of a $R$-regular sequence $x_1, \ldots, x_d$ of the maximal
ideal such that $R/(x_1, \ldots, x_d)$ has dimension $0$.
We will usually just say ``regular sequence'' and not
``$R$-regular sequence''.

\begin{lemma}
\label{lemma-reformulate-CM}
\begin{slogan}
Regular sequences in Cohen-Macaulay local rings are characterized
by cutting out something of the correct dimension.
\end{slogan}
Let $R$ be a Noetherian local Cohen-Macaulay ring with maximal
ideal $\mathfrak m $. Let $x_1, \ldots, x_c \in \mathfrak m$ be
elements. Then
$$
x_1, \ldots, x_c
\text{ is a regular sequence }
\Leftrightarrow
\dim(R/(x_1, \ldots, x_c)) = \dim(R) - c
$$
If so
$x_1, \ldots, x_c$ can be extended to
a regular sequence of length $\dim(R)$ and each quotient
$R/(x_1, \ldots, x_i)$ is a Cohen-Macaulay ring of dimension
$\dim(R) - i$.
\end{lemma}

\begin{proof}
Special case of Proposition \ref{proposition-CM-module}.
\end{proof}

\begin{lemma}
\label{lemma-maximal-chain-CM}
Let $R$ be Noetherian local.
Suppose $R$ is Cohen-Macaulay of dimension $d$.
Any maximal chain of ideals $\mathfrak p_0 \subset
\mathfrak p_1 \subset \ldots \subset \mathfrak p_n$
has length $n = d$.
\end{lemma}

\begin{proof}
Special case of Lemma \ref{lemma-maximal-chain-maximal-CM}.
\end{proof}

\begin{lemma}
\label{lemma-CM-dim-formula}
Suppose $R$ is a Noetherian local Cohen-Macaulay ring of dimension $d$.
For any prime $\mathfrak p \subset R$ we have
$$
\dim(R) = \dim(R_{\mathfrak p}) + \dim(R/\mathfrak p).
$$
\end{lemma}

\begin{proof}
Follows immediately from Lemma \ref{lemma-maximal-chain-CM}.
(Also, this is a special case of Lemma \ref{lemma-dim-formula-maximal-CM}.)
\end{proof}

\begin{lemma}
\label{lemma-localize-CM}
Suppose $R$ is a Cohen-Macaulay local ring.
For any prime $\mathfrak p \subset R$ the
ring $R_{\mathfrak p}$ is Cohen-Macaulay as well.
\end{lemma}

\begin{proof}
Special case of Lemma \ref{lemma-localize-CM-module}.
\end{proof}

\begin{definition}
\label{definition-ring-CM}
A Noetherian ring $R$ is called {\it Cohen-Macaulay} if all
its local rings are Cohen-Macaulay.
\end{definition}

\begin{lemma}
\label{lemma-CM-polynomial-algebra}
Suppose $R$ is a Noetherian Cohen-Macaulay ring.
Any polynomial algebra over $R$ is Cohen-Macaulay.
\end{lemma}

\begin{proof}
Special case of Lemma \ref{lemma-maximal-CM-polynomial-algebra}.
\end{proof}

\begin{lemma}
\label{lemma-dimension-shift}
Let $R$ be a Noetherian local Cohen-Macaulay ring of dimension $d$.
Let $0 \to K \to R^{\oplus n} \to M \to 0$ be an exact sequence of
$R$-modules. Then either $M = 0$, or $\text{depth}(K) > \text{depth}(M)$, or
$\text{depth}(K) = \text{depth}(M) = d$.
\end{lemma}

\begin{proof}
This is a special case of Lemma \ref{lemma-depth-in-ses}.
\end{proof}

\begin{lemma}
\label{lemma-mcm-resolution}
Let $R$ be a local Noetherian Cohen-Macaulay ring of dimension $d$.
Let $M$ be a finite $R$-module of depth $e$.
There exists an exact complex
$$
0 \to K \to F_{d-e-1} \to \ldots \to F_0 \to M \to 0
$$
with each $F_i$ finite free and $K$ maximal Cohen-Macaulay.
\end{lemma}

\begin{proof}
Immediate from the definition and Lemma \ref{lemma-dimension-shift}.
\end{proof}

\begin{lemma}
\label{lemma-find-sequence-image-regular}
Let $\varphi : A \to B$ be a map of local rings.
Assume that $B$ is Noetherian and Cohen-Macaulay and that
$\mathfrak m_B = \sqrt{\varphi(\mathfrak m_A) B}$. Then there exists
a sequence of elements $f_1, \ldots, f_{\dim(B)}$ in $A$
such that $\varphi(f_1), \ldots, \varphi(f_{\dim(B)})$ is a
regular sequence in $B$.
\end{lemma}

\begin{proof}
By induction on $\dim(B)$ it suffices to prove: If $\dim(B) \geq 1$, then we
can find an element $f$ of $A$ which maps to a nonzerodivisor in $B$.
By
Lemma \ref{lemma-reformulate-CM}
it suffices to find $f \in A$ whose image in $B$ is not contained in any
of the finitely many minimal primes $\mathfrak q_1, \ldots, \mathfrak q_r$
of $B$. By the assumption that
$\mathfrak m_B = \sqrt{\varphi(\mathfrak m_A) B}$
we see that $\mathfrak m_A \not \subset \varphi^{-1}(\mathfrak q_i)$.
Hence we can find $f$ by
Lemma \ref{lemma-silly}.
\end{proof}








